% R568: fuente francesa intro-fr.tex:95--132, pp. 8--9.
Por último, huelga decir que un libro de geometría algebraica, y sobre todo uno dedicado a sus fundamentos, está necesariamente influido, aunque sea a través de otras personas, por matemáticos como O.~Zariski y A.~Weil.
En particular, la \emph{Théorie des fonctions holomorphes} de Zariski~\cite{I-20}, convenientemente flexibilizada mediante métodos cohomológicos y completada con un teorema de existencia (cap.~III, \S\S~4 y 5), es, junto con la técnica de descenso expuesta en el cap.~VI, una de las principales herramientas empleadas en este tratado, y nos parece una de las más potentes de que dispone la geometría algebraica.

La técnica general en la que se inserta puede esbozarse del siguiente modo (se dará un ejemplo típico en el cap.~IX, al estudiar el grupo fundamental).
Se tiene un morfismo propio (cap.~II) $f:X\to Y$ de una variedad algebraica en otra (más generalmente, de un esquema en otro), que se quiere estudiar cerca de un punto $y\in Y$ para resolver un problema P relativo a una vecindad de $y$.
Se procede por etapas sucesivas:
\begin{enumerate}
  \item[$1^\circ$] Puede suponerse que $Y$ es afín, de modo que $X$ pasa a ser un esquema definido sobre el anillo afín $A$ de $Y$; incluso puede sustituirse $A$ por el anillo local de $y$.
  Esta reducción siempre es fácil en la práctica (cap.~V) y nos lleva al caso en que $A$ es un anillo \emph{local}.
  \item[$2^\circ$] Se estudia el problema planteado cuando $A$ es un anillo local \emph{artiniano}.
  Para que siga teniendo efectivamente sentido cuando no se supone que $A$ sea un dominio, a veces hay que reformular el problema P; de este modo, a menudo se obtiene una mejor comprensión del problema, de carácter «infinitesimal» en esta etapa.
  \item[$3^\circ$] La teoría de esquemas formales (cap.~III, \S\S~3, 4 y 5) permite pasar del caso de un anillo artiniano al de un \emph{anillo local completo}.
  \item[$4^\circ$] Por último, si $A$ es un anillo local cualquiera, la consideración de «secciones multiformes» sobre esquemas adecuados sobre $X$, que aproximan una sección «formal» dada (cap.~IV), permitirá a menudo pasar de un resultado conocido para el esquema

  \oldpage[I]{9}
  deducido de $X$ por extensión de escalares al completado de $A$, a un resultado análogo para una extensión finita bastante sencilla (por ejemplo, no ramificada) de $A$.
\end{enumerate}

Este esbozo muestra la importancia del estudio sistemático de los esquemas definidos sobre un anillo artiniano $A$.
El punto de vista de Serre en su formulación de la teoría local de cuerpos de clases y algunos trabajos recientes de Greenberg parecen sugerir que ese estudio podría emprenderse asociando funtorialmente a tal esquema $X$ un esquema $X'$ sobre el cuerpo residual $k$ de $A$ (suponiendo perfecto ese cuerpo), cuya dimensión sea igual, en los casos favorables, a $n\dim X$, donde $n$ es la longitud de $A$.

En cuanto a la influencia de A.~Weil, baste decir que la necesidad de desarrollar las herramientas para formular con toda la generalidad deseada la definición de «cohomología de Weil» y abordar la demostración\footnote{Para evitar cualquier malentendido, precisemos que esta tarea apenas se ha iniciado cuando se escribe esta introducción y, por tanto, todavía no ha dado lugar a una demostración de las conjeturas de Weil.} de todas las propiedades formales necesarias para establecer sus célebres conjeturas en geometría diofántica~\cite{I-19} ha sido una de las principales motivaciones para redactar este tratado, al igual que el deseo de encontrar el marco natural de las nociones y los métodos habituales de la geometría algebraica y dar a los autores la ocasión de comprender tales nociones y técnicas.

\bigskip
\begin{center}*\quad*\quad*\end{center}
\bigskip

Para terminar, creemos útil advertir a los lectores que, al igual que los propios autores, sin duda encontrarán ciertas dificultades antes de acostumbrarse al lenguaje de los esquemas y convencerse de que las construcciones habituales sugeridas por la intuición geométrica pueden transcribirse en ese lenguaje, esencialmente de una sola manera razonable.
Como en muchas partes de la matemática moderna, la intuición inicial parece alejarse cada vez más del lenguaje apropiado para expresarla con toda la precisión y generalidad deseadas.
En este caso, la dificultad psicológica consiste en la necesidad de trasladar a los objetos de una categoría ya bastante distinta de la categoría de conjuntos (a saber, la categoría de preesquemas o la de preesquemas sobre un preesquema dado) nociones familiares para los conjuntos: productos cartesianos, leyes de grupo, de anillo y de módulo, fibrados, fibrados principales homogéneos, etc.
En el futuro, sin duda será difícil para el matemático eludir este nuevo esfuerzo de abstracción, quizá bastante pequeño, en definitiva, comparado con el que realizaron nuestros predecesores al familiarizarse con la teoría de conjuntos.

\bigskip
\begin{center}*\quad*\quad*\end{center}
\bigskip

Las referencias se darán siguiendo el sistema decimal; por ejemplo, en III, 4.9.3, III indica el capítulo, el número~4 el parágrafo y el~9 la sección de ese parágrafo.
Dentro de un mismo capítulo se omitirá la indicación del capítulo.
